\section{Background}
\label{sec:background}

% =========================================================================
% RL-first pedagogical ordering, per design.md section 1.6:
%   MDP -> Bellman -> Q-learning -> approximate Q-learning -> exploration
%   -> MCTS as sample-efficient policy improvement -> Expert Iteration
%   -> LLM deliberation as a context MDP.
% Style note (from design.md section 12): expanded sums before expectation
% notation; define every symbol before it is abbreviated.
% =========================================================================

This section assembles the reinforcement-learning vocabulary used in the
rest of the paper. We move from the Markov decision process and Bellman
optimality through tabular and approximate $Q$-learning to MCTS, framed
as a procedure that supplies stronger training targets to the same
learning rule. Expert iteration and the LLM-as-environment framing then
follow.

\subsection{Markov Decision Processes}
\label{sec:background-mdp}

A finite discounted Markov decision process is a tuple
$\Mset = (\Sset, \Aset, P, R, \gamma)$, where $\Sset$ is a state space,
$\Aset(s) \subseteq \Aset$ is the set of actions available in state $s$,
$P(s' \given s, a)$ is the probability of transitioning to state $s'$ after
taking action $a$ in state $s$, $R(s, a, s') \in \R$ is the immediate reward
for that transition, and $\gamma \in [0, 1)$ is the discount factor.

The Markov assumption states that, conditional on $s$ and $a$, the next
state $s'$ is independent of the prior history,
\begin{equation}
P(s_{t+1} \given s_0, a_0, \ldots, s_t, a_t) = P(s_{t+1} \given s_t, a_t).
\label{eq:markov}
\end{equation}
When history matters, the state representation must incorporate the
relevant history; the next section discusses one such case for LLM
deliberation, where the practical ``state'' is a structured context that
may only approximate a Markov state.

We will care about two computational presentations of an MDP. A
\emph{finite MDP} exposes the transition kernel $P$ and reward function
$R$ explicitly as tables, allowing exact enumeration of $\sum_{s'}$ in the
Bellman equation. A \emph{generative MDP} exposes only a sampler: given
$(s, a)$, it returns a single draw $s' \sim P(\cdot \given s, a)$ and the
corresponding reward $r = R(s, a, s')$. The generative interface is what
remains when $P$ is too large to tabulate, when it is unknown, or when it
is implicitly defined by a process such as an LLM forward pass.
Section~\ref{sec:method-mcts} formalizes this distinction as the
\texttt{FiniteMDP} and \texttt{GenerativeMDP} protocols and explains why
MCTS can be built directly on the latter while value iteration requires
the former.

\subsection{Bellman Optimality and the Action-Value Function}
\label{sec:background-bellman}

The optimal state-value function $\Vstar$ and the optimal action-value
function $\Qstar$ satisfy the Bellman optimality equations,
\begin{equation}
\Qstar(s, a)
= \sum_{s'} P(s' \given s, a)
  \left[ R(s, a, s') + \gamma \max_{a'} \Qstar(s', a') \right],
\label{eq:bellman}
\end{equation}
\begin{equation}
\Vstar(s) = \max_{a \in \Aset(s)} \Qstar(s, a).
\label{eq:vstar}
\end{equation}
The right-hand side of \eqref{eq:bellman} decomposes into three steps:
(i)~enumerate the possible next states $s'$ weighted by their transition
probabilities $P(s' \given s, a)$, (ii)~for each transition collect the
immediate reward $R(s, a, s')$, and (iii)~add the discounted optimal
future value $\gamma \max_{a'} \Qstar(s', a')$. Both equations make
explicit the sum that any sample-based method must eventually approximate.

The greedy deterministic policy derived from $\Qstar$ is
\begin{equation}
\pi^{\ast}(s) = \arg\max_{a \in \Aset(s)} \Qstar(s, a),
\label{eq:greedy-policy}
\end{equation}
and a smooth stochastic alternative is the softmax policy with temperature
$\tau_{\pi} > 0$,
\begin{equation}
\pi_{\tau_{\pi}}(a \given s)
= \frac{\exp(\Qstar(s, a) / \tau_{\pi})}{\sum_{b \in \Aset(s)} \exp(\Qstar(s, b) / \tau_{\pi})}.
\label{eq:softmax-policy}
\end{equation}
As $\tau_{\pi} \to 0$, $\pi_{\tau_{\pi}}$ concentrates on $\pi^{\ast}$; as
$\tau_{\pi} \to \infty$, it approaches the uniform distribution over $\Aset(s)$.
Soft policies are convenient for exploration and for differentiable
training, both of which appear later.

\subsection{Q-Learning from Sampled Transitions}
\label{sec:background-qlearning}

When $P$ and $R$ are known and $\Sset$ is small enough to tabulate, the
sum in \eqref{eq:bellman} can be evaluated exactly via value iteration. In
the reinforcement-learning setting we usually cannot. Instead, the agent
observes a single sampled transition
\begin{equation}
(s, a, r, s'), \qquad s' \sim P(\cdot \given s, a), \qquad r = R(s, a, s'),
\label{eq:sample-transition}
\end{equation}
and updates its value estimate from that one sample. The tabular
$Q$-learning update of \citet{sutton2018reinforcement} is
\begin{equation}
Q(s, a) \;\gets\; (1 - \alpha)\, Q(s, a)
                + \alpha \left[ r + \gamma \max_{a'} Q(s', a') \right],
\label{eq:qlearning}
\end{equation}
where $\alpha \in (0, 1]$ is a step size. Equivalently,
\begin{equation}
Q(s, a) \;\gets\; Q(s, a)
              + \alpha \left( \left[ r + \gamma \max_{a'} Q(s', a') \right]
                              - Q(s, a) \right),
\label{eq:qlearning-error}
\end{equation}
which exposes the Bellman temporal-difference error in brackets.

The quantity
\begin{equation}
y = r + \gamma \max_{a'} Q(s', a')
\label{eq:bellman-target}
\end{equation}
is the \emph{Bellman target} computed from a single sampled transition. It
should be contrasted with the full inner expression on the right-hand
side of \eqref{eq:bellman}: the target $y$ uses one realization of $s'$
in place of the explicit sum $\sum_{s'} P(s' \given s, a) [\cdots]$.
Averaged over many samples drawn from $P(\cdot \given s, a)$, the
expected target equals the right-hand side of \eqref{eq:bellman}, which
is what makes the update a stochastic approximation of the Bellman
operator.

A key property of \eqref{eq:qlearning} is that it is \emph{off-policy}.
The bootstrap term $\max_{a'} Q(s', a')$ refers to the greedy continuation,
not to the action that any particular behavior policy would have taken at
$s'$. As a consequence, $Q$-learning can be trained from experience
generated by any behavior policy $\mu$ that gives sufficient coverage of
state-action pairs, including a uniform random policy, an
$\epsilon$-greedy policy, a softmax policy over the current $Q$, or a
search procedure such as MCTS. This decoupling between the behavior
policy that samples experience and the target policy that the updates
converge toward is what makes MCTS-guided sample collection (developed
later in this section) compatible with $Q$-learning at all.

Under standard conditions (Robbins-Monro step sizes
$\sum_t \alpha_t = \infty$, $\sum_t \alpha_t^2 < \infty$, and a behavior
policy that visits every state-action pair infinitely often),
\eqref{eq:qlearning} converges almost surely to $\Qstar$
\citep{sutton2018reinforcement}. These conditions never hold exactly in
practice; the conceptual picture of sample-based stochastic
approximation of the Bellman operator remains the right one.

\subsection{Approximate Q-Learning}
\label{sec:background-approxq}

Tabular $Q$-learning stores one scalar per state-action pair and is
infeasible when $\Sset$ is large, continuous, or textual. In the
LLM-deliberation setting of this paper the state is a structured context
of tokens, so a table is out of the question. We replace the table with a
parametric action-value function $\Qpsi(s, a)$ with parameters $\psi$.

For a sampled transition $(s, a, r, s')$, the Bellman target is now
\begin{equation}
y = r + \gamma \max_{a'} Q_{\bar{\psi}}(s', a'),
\label{eq:bellman-target-approx}
\end{equation}
where $Q_{\bar{\psi}}$ is a \emph{target network}: a slow-moving copy of
$\Qpsi$ whose parameters $\bar{\psi}$ are periodically synchronized from
$\psi$ \citep{mnih2015human}. The parameters $\psi$ are then updated by
gradient descent on the squared Bellman error,
\begin{equation}
\mathcal{L}_Q(\psi) \;=\; \tfrac{1}{2}\,
\E_{(s, a, r, s') \sim \mathcal{D}}
\left[ \left( \Qpsi(s, a) - y \right)^2 \right],
\label{eq:approx-q-loss}
\end{equation}
where the expectation is over transitions drawn from a dataset
$\mathcal{D}$ of past experience. In practice $\mathcal{D}$ is a
\emph{replay buffer}, a finite ring of recent transitions sampled
uniformly or with priority; the buffer breaks the strong temporal
correlation between consecutive transitions in a trajectory.

Practical approximate $Q$-learning bundles together a small collection of
stabilization techniques: the replay buffer just described, the target
network in \eqref{eq:bellman-target-approx}, clipped or normalized
rewards to keep the regression target bounded, exploration schedules that
anneal a behavior temperature or $\epsilon$ over training, and various
forms of regularization on $\psi$. None of these change the core idea,
which is simply
\begin{equation}
\text{sample experience} \;\longrightarrow\;
\text{compute Bellman target}~y \;\longrightarrow\;
\text{minimize } (\Qpsi(s, a) - y)^2.
\label{eq:approxq-loop}
\end{equation}
Section~\ref{sec:method-architecture} instantiates $\Qpsi$ as a small
action-value head over a frozen LLM backbone, with LoRA adaptation
\citep{hu2022lora} of the backbone reserved as a phase-2 option.

\subsection{Exploration and Behavior Policies}
\label{sec:background-exploration}

Equations \eqref{eq:bellman-target} and \eqref{eq:bellman-target-approx}
say nothing about how the transitions $(s, a, r, s')$ feeding the
$Q$-update are generated. That choice is the responsibility of the
\emph{behavior policy} $\mu(a \given s)$, and it interacts with learning
through coverage: any state-action pair that $\mu$ never visits cannot be
updated.

Common behavior policies include the \emph{uniform random} policy
$\mu(a \given s) = 1 / |\Aset(s)|$ (maximally exploratory, useful as a
warm start); the \emph{greedy} policy that always picks
$\arg\max_a \Qpsi(s, a)$ (exploits current knowledge but can collapse
onto a suboptimal action under biased $\Qpsi$); the
\emph{$\epsilon$-greedy} policy that flips between random and greedy
with probability $\epsilon$ (typically annealed over training); the
\emph{softmax / Boltzmann} policy
$\mu(a \given s) \propto \exp(\Qpsi(s, a) / \tau_{\pi})$ with temperature
$\tau_{\pi}$; and the \emph{search-guided} policy
$\mu(a \given s) = \pimcts(a \given s)$ produced by an online procedure
such as MCTS that uses $\Qpsi$ and a policy prior $\ptheta$ to focus
sampling on promising actions.

The classical exploration-exploitation tradeoff
\citep{sutton2018reinforcement} appears as a tension between a behavior
policy that exploits too greedily, missing high-return regions, and one
that explores too aggressively, wasting samples on actions already known
to be poor. Search-guided behavior is attractive because it adapts the
exploration distribution to local context using current value and policy
estimates, rather than treating exploration as a state-independent
schedule.

\subsection{MCTS as Sample-Efficient Policy Improvement}
\label{sec:background-mcts}

Monte Carlo tree search, in the form introduced by
\citet{kocsis2006bandit} and \citet{coulom2006efficient} and surveyed by
\citet{browne2012survey}, is best understood here not as a standalone
planner but as a procedure that, given a current parametric estimate
$(\Qpsi, \ptheta)$, produces \emph{locally improved} value and policy
estimates $(\Qmcts, \pimcts)$ at a query state $s_0$. Those improved
estimates then serve as targets for further training of $(\Qpsi,
\ptheta)$ via the same approximate $Q$-learning machinery from
section~\ref{sec:background-approxq}.

\paragraph{Sampling replaces summation.}
Equation~\eqref{eq:bellman} contains the sum
$\sum_{s'} P(s' \given s, a) [\cdots]$. When $\Sset$ is large or $P$ is
implicit, this sum is intractable. MCTS replaces it with samples: at any
$(s, a)$ pair encountered in the tree, the search calls the generative
MDP for a draw $s' \sim P(\cdot \given s, a)$ and propagates the
resulting reward and next-state value back along the path. The empirical
chance-node mean
\begin{equation}
\Qmcts(s, a) \;=\; \frac{W(s, a)}{N(s, a)}
\;=\; \frac{1}{N(s, a)} \sum_{i=1}^{N(s, a)} G_i,
\label{eq:mcts-q-estimate}
\end{equation}
where $G_i$ is the discounted return observed on the $i$th simulation
passing through $(s, a)$, $W(s, a)$ is the cumulative return, and
$N(s, a)$ is the visit count, is by construction a sample estimator of
the Bellman expectation in \eqref{eq:bellman}. As $N(s, a) \to \infty$
with consistent rollout values, $\Qmcts(s, a) \to \Qstar(s, a)$.

\paragraph{Chance-node bifurcation.}
Because LLM-driven transitions are stochastic, the tree distinguishes
\emph{decision nodes} (states $s$ where the controller chooses an
action) from \emph{chance nodes} ($(s, a)$ pairs where the environment
draws $s'$). The counts $N(s, a)$ and returns $W(s, a)$ are stored at
chance nodes; each distinct sampled successor becomes a child decision
node. This bifurcation makes \eqref{eq:mcts-q-estimate} an honest
estimator of the Bellman expectation rather than a degenerate trace
through a single imagined trajectory.

\paragraph{Per-edge reward backup.}
Returns are computed backward along the path of (chance node, sampled
reward) pairs visited during a simulation. If a rollout traverses the
sequence of $(s_t, a_t, r_{t+1}, s_{t+1})$ for $t = 0, 1, \ldots, T-1$ and
terminates with a leaf value $\hat V(s_T)$, the return at depth $t$ is
\begin{equation}
G_t = r_{t+1} + \gamma\, G_{t+1}, \qquad G_T = \hat V(s_T),
\label{eq:per-edge-backup}
\end{equation}
and $G_t$ is added to the running statistics $W(s_t, a_t)$ and
$N(s_t, a_t)$ of every ancestor chance node. The crucial point, often
missed in tutorial presentations that only credit terminal rewards, is
that each intermediate reward $r_{t+1}$ enters the backup at its own
discount level. Section~\ref{sec:method-mcts} works through this in
implementation detail.

\paragraph{PUCT selection.}
At each decision node $s$ in the tree, the search selects an action by
maximizing the PUCT score \citep{silver2017alphagozero},
\begin{equation}
\PUCT(s, a) \;=\; \widehat Q(s, a)
\;+\; c_{\mathrm{puct}} \;\ptheta(a \given s)\;
\frac{\sqrt{N(s)}}{1 + N(s, a)},
\label{eq:puct-bg}
\end{equation}
where $\widehat Q(s, a)$ is the current empirical estimate at the
chance node (or the parametric prior $\Qpsi(s, a)$ if the node is
unvisited), $\ptheta(a \given s)$ is the learned policy prior,
$N(s) = \sum_{a'} N(s, a')$, and $c_{\mathrm{puct}} > 0$ is an
exploration constant. The first term encourages exploitation of actions
that already look good; the second encourages exploration of actions
with high prior probability that have not yet been sampled enough.

\paragraph{Visit-count policy.}
After a fixed simulation budget at root $s_0$, the search returns the
visit-count distribution
\begin{equation}
\pimcts(a \given s_0)
\;=\; \frac{N(s_0, a)^{1/\tau}}
            {\sum_{a' \in \Aset(s_0)} N(s_0, a')^{1/\tau}},
\label{eq:visit-policy}
\end{equation}
where $\tau > 0$ controls sharpness. Under reasonable conditions, this
visit-count policy is a stronger policy than the input prior $\ptheta$:
PUCT preferentially expands actions that have either high prior or high
empirical $Q$, so $\pimcts$ concentrates on actions that look good both
under the prior and under the additional information accumulated by
search.

\paragraph{Sample-efficient policy improvement.}
MCTS thus plays two roles inside an RL loop. It is a sample-efficient
\emph{behavior policy} whose root actions are typically much better than
those of the prior or current $Q$ alone, and it produces
\emph{stronger training targets}: the chance-node aggregates
$\Qmcts(s, a)$ average over many simulated continuations and are
lower-variance than any single rollout return. Approximate $Q$-learning
consumes these targets exactly as it would any other Bellman targets,
by minimizing the squared error in \eqref{eq:approx-q-loss}.

\subsection{Expert Iteration as Joint Q and Policy Distillation}
\label{sec:background-expit}

Expert iteration \citep{anthony2017thinking} combines the components of
the preceding subsections into a single training loop. AlphaGo Zero and
AlphaZero \citep{silver2017alphagozero, silver2018alphazero} are the
best-known realizations: in those systems, the parametric policy and
value networks guide MCTS, MCTS improves on them locally, and the
improvement is distilled back into the parameters. The structure we use
in this paper is the same, expressed in the RL vocabulary of the
preceding subsections.

The loop iterates the following steps.
\begin{enumerate}[leftmargin=*,itemsep=2pt]
  \item \textbf{Sample from a search-guided behavior policy.} At each
    encountered state $s$, run MCTS with $\Qpsi$ as the value prior at
    leaves and $\ptheta$ as the policy prior in
    \eqref{eq:puct-bg}. Collect the chance-node statistics
    $\Qmcts(s, a)$ from \eqref{eq:mcts-q-estimate} and the root
    visit-count distribution $\pimcts(\cdot \given s)$ from
    \eqref{eq:visit-policy}. Sample the next environment action from
    $\pimcts$, generating a transition $(s, a, r, s')$.
  \item \textbf{Distill the value into $\Qpsi$.} Minimize the squared
    Bellman error of \eqref{eq:approx-q-loss}, using $\Qmcts(s, a)$ as
    the target,
    \begin{equation}
    \mathcal{L}_Q(\psi)
    \;=\; \tfrac{1}{2}\, \E \left[ (\Qpsi(s, a) - \Qmcts(s, a))^2 \right].
    \label{eq:expit-q-loss}
    \end{equation}
    This is approximate $Q$-learning, with MCTS supplying both the
    sampling distribution and the targets.
  \item \textbf{Distill the policy into $\ptheta$.} Minimize the
    cross-entropy of the visit-count policy under the policy prior,
    \begin{equation}
    \mathcal{L}_{\pi}(\theta)
    \;=\; -\,\E_s \left[
      \sum_{a \in \Aset(s)} \pimcts(a \given s) \log \ptheta(a \given s)
    \right].
    \label{eq:expit-pi-loss}
    \end{equation}
  \item \textbf{Iterate.} The improved $(\Qpsi, \ptheta)$ become the
    priors for the next round of search.
\end{enumerate}

The two distillation losses share a common structure. Each one trains a
parametric model toward a target produced by a stronger but more
expensive procedure (search) that itself used the model as input. This
is what \citet{anthony2017thinking} call ``thinking, fast and slow'':
the parametric models are the fast system, MCTS is the slow system, and
training shifts competence from the slow system into the fast one.

\subsection{LLM Deliberation as a Context MDP}
\label{sec:background-llm-mdp}

The framing above is generic. To apply it to language-model
deliberation we now specify how an LLM induces the MDP, following the
design of section~\ref{sec:method-state}.

\paragraph{The LLM as environment.}
An LLM together with a fixed prompt template induces a conditional
distribution over next contexts given a control action,
\begin{equation}
C_{t+1} \;\sim\; P_{\mathrm{LLM}}(\cdot \given C_t, a_t),
\label{eq:llm-transition}
\end{equation}
where $C_t$ is a structured context (the running summary, candidate
solutions, critiques, verifications, and depth produced so far) and
$a_t$ is a control action such as ``propose a solution'', ``critique
the current candidate'', or ``verify the current candidate''.
Section~\ref{sec:method-actions} fixes the action vocabulary. The
controller's experience is exactly the sequence of sampled transitions
$(C_t, a_t, r_t, C_{t+1})$.

\paragraph{Model-free from the controller's perspective.}
Although \eqref{eq:llm-transition} is, in principle, a transition model,
the controller does not enumerate it. It observes sampled transitions
from the LLM and updates $\Qpsi(C, a)$ using
\eqref{eq:bellman-target-approx} and \eqref{eq:approx-q-loss}. From the
controller's perspective this is model-free RL: the LLM is the
environment interface, not a separate simulator of an external world.
MCTS, when used here, draws samples from the same LLM environment
according to \eqref{eq:llm-transition}, organized as search.

\paragraph{Partial observability and the Markov bet.}
The context $C$ is a summary of dialogue and intermediate reasoning, not
a sufficient statistic of the full history of LLM token-level
realizations. Strictly speaking the process is a partially observed
MDP, and treating $C$ as a state is a state-abstraction approximation.
This is the standard situation in deep RL, where the input to the value
function is an observation or learned representation rather than a true
Markov state. The practical bet, encoded in our choice of context
representation in section~\ref{sec:method-state}, is that the
state-abstraction approximation is good enough for $\Qpsi(C, a)$ to be
useful, and that the MCTS chance-node statistics in
\eqref{eq:mcts-q-estimate} still concentrate on values that are useful
for action selection even when the underlying process is not exactly
Markov in $C$.
